\documentclass[10pt,a4paper]{amsart}
\usepackage[margin=19mm]{geometry}
\usepackage{amsmath,amssymb,mathtools}
\usepackage[scr=rsfs,cal=euler]{mathalfa}
\usepackage{xfrac}
\usepackage[unicode,hidelinks,bookmarks=false]{hyperref}
\newcommand{\ie}{i.e.\ }
\newcommand{\cf}{cf.\ }
\newcommand{\resp}{respectively\ }
\newcommand{\Subsection}{\subsection}
\providecommand{\nobreakdash}{\nobreak\hbox{-}}
\setlength{\emergencystretch}{2em}
\multlinegap0pt
\usepackage{bm}
\usepackage{bbm}
\usepackage{dsfont}
\usepackage{xspace}
\usepackage{cancel}
\usepackage{mathtools}
\usepackage{amsthm}
\makeatletter
\@ifundefined{lemma}{\newtheorem{lemma}{Lemma}}{}
\makeatother
\newcommand{\semantics}[1]{[\![\mbox{\em $ #1 $\/}]\!]}
\title[Interleaving Logic and Counting]{Interleaving Logic and Counting}
\author{Johan van Benthem, Thomas Icard}
\date{Source: arXiv:2507.05219v1 (CC BY 4.0)}
\begin{document}
\maketitle

\noindent\emph{Source and licence.} Interleaving Logic and Counting. Version arXiv:2507.05219v1 (CC BY 4.0), distributed under the Creative Commons Attribution 4.0 International licence, as recorded in the arXiv licence metadata for this exact version and shown on the abstract page https://arxiv.org/abs/2507.05219v1. Full rights notice: licenses/arxiv-2507.05219-CC-BY-4.0.txt.\par\smallskip

\noindent\textit{Editorial scope.} Counted unit: the Lemma whose source label is \texttt{None} in the article (source0.tex lines 693--694, source label \texttt{None}), delivered together with the article's own complete original proof of it (source0.tex lines 695--722, one \texttt{proof} environment of 28 lines). No mathematics is written, repaired or re-derived: statement and proof are the article's own text line for line. Required context retained uncounted: none. Every definition, display and label of those lines is retained unchanged. Omitted from this package and not redistributed: 1--692, 723--2368. Nothing inside a retained excerpt is omitted or altered: every retained body is delivered exactly as the article's lines are written, so no omission receipt of any kind accompanies this package. Citations: the retained excerpts carry no citation command at all, so no bibliography entry of the article is transcribed and no reference is redistributed. Reference adaptation: none was needed and none was made; every reference inside the delivered unit resolves inside it, so no unresolved reference or citation is produced. Printed slips kept literal: no printed word, symbol, hypothesis or number of the counted statement or of its proof was altered. Layout and class: the article's own class is \texttt{amsart} with packages including inputenc, amsmath, amssymb, latexsym, natbib, amsthm, amsfonts, fullpage, color, graphicx, microtype, hyperref, enumitem, tikz; the wrapper is the lane's pinned \texttt{amsart} editorial wrapper at 10pt on A4, which restates the theorem-like environments the retained text needs and adds the article's own macro definitions that the retained text needs (\texttt{\textbackslash semantics}), with extra wrapper packages (bm, bbm, dsfont, xspace, cancel, mathtools). The delivered numbering is the unit's own. Assets: no retained span contains a figure environment, a graphics-inclusion command, a PDF-page inclusion, a table or a TikZ picture; no image file is copied into this package, the package declares no asset dependency and no image file is redistributed with it. Licence: version arXiv:2507.05219v1 of this article is distributed under the Creative Commons Attribution 4.0 International licence, as recorded in the arXiv licence metadata for this exact version and shown on the abstract page https://arxiv.org/abs/2507.05219v1. A full rights notice is delivered at licenses/arxiv-2507.05219-CC-BY-4.0.txt. Characters: every retained excerpt is pure ASCII, so the delivered engine, XeLaTeX with its default Latin Modern text font, reports no missing character.
\begin{lemma} Let $X$ be a predicate variable, and $O_1,\dots,O_n$ be either predicate letters or predicate variables. Suppose that $\varphi(O_1,\dots,O_n,X)$ defines a linear set in state-descriptions $S_1,\dots,S_{2^{n+1}}$ over $O_1,\dots,O_n,X$. Then $\exists X. \varphi(O_1,\dots,O_n,X)$ defines a linear set in state-descriptions $S_1',\dots,S_{2^n}'$ over just $O_1,\dots,O_n$.
\end{lemma}

\begin{proof} As $\varphi(O_1,\dots,O_n,X)$ is linear, we can assume it defines the solutions to  \vspace{0.5ex} \begin{equation}
 \begin{pmatrix}
\;|S_1|\;  \\
\vdots  \\
|S_{2^{n+1}}|  
\end{pmatrix} \hspace{.1in}= \hspace{.1in}
\begin{pmatrix}
b_1 + a_{1,1}\mathsf{u}_1 + & \dots  & + a_{1,m}\mathsf{u}_m  \\
 \vdots & & \vdots \\
b_{2^{n+1}} + a_{2^{n+1},1}\mathsf{u}_1 + & \dots & +  a_{2^{n+1},m}\mathsf{u}_m
\end{pmatrix}. \label{semi-lin}
\end{equation}\vspace{0.5ex} To show that $\exists X. \varphi(O_1,\dots,O_n,X)$, too, is linear, we define another linear set of equations by ``projecting out'' the variable $X$. 

\vspace{1ex}

Specifically, note for each state-description $S_k'$, both $S_k' \wedge X$ and $S_k' \wedge \neg X$ are (equivalent to) some state-descriptions, $S_i$ and $S_j$, and in fact $S_k'$ is equivalent to $S_i \vee S_j$. The new linear system in $2^n$ variables is then as follows for each $k\leq 2^n$: \begin{eqnarray} |S'_k| & = &  b_i + b_j + (a_{i,1}+a_{j,1})\mathsf{u}_1+ \dots + (a_{i,m} + a_{j,m})\mathsf{u}_m . \label{ex-line}\end{eqnarray} It remains only to show that:\vspace{0.5ex} \begin{eqnarray}
   \mathcal{M} \vDash \exists X. \varphi(O_1,\dots,O_n,X) 
   & \Leftrightarrow & [|S_1'|_{\mathcal{M}},\dots,|S_{2^n}'|_{\mathcal{M}}] \mbox{ is a solution to equations }(\ref{ex-line}).\vspace{1ex} \nonumber
\end{eqnarray} 

$(\Rightarrow)$: If $\mathcal{M} \vDash \exists X. \varphi(O_1,\dots,O_n,X)$ then for some subset $A$ of the domain, we have $\mathcal{M},s^X_A \vDash \varphi(O_1,\dots,O_n,X)$. Treating $X$ now as a predicate constant, we have a model $\mathcal{M}'$ for which $X^{\mathcal{M}'}=A$, and by assumption this gives a solution $[|S_1|_{\mathcal{M}'},\dots,|S_{2^{n+1}}|_{\mathcal{M}'}]$ to (\ref{semi-lin}). But each state-description $S_k'$ is equivalent to a disjunction $S_i \vee S_j$, whose cardinality is the sum $|S_i|+|S_j|$. Therefore $\mathcal{M}'$ will satisfy each of the constraints in (\ref{ex-line}). As the state-descriptions $S_1',\dots,S_{2^n}'$ are independent of $X$, this means  $[|S_1'|_{\mathcal{M}},\dots,|S_{2^n}'|_{\mathcal{M}}]$ also gives a solution to (\ref{ex-line}).

\vspace{1ex}

$(\Leftarrow)$: Suppose $[|S_1'|_{\mathcal{M}},\dots,|S_{2^n}'|_{\mathcal{M}}]$  gives a solution to (\ref{ex-line}) for some particular choices $\mathsf{u}_1,\dots,\mathsf{u}_m$. We need to find a set $A$ such that $\mathcal{M},s^X_A \vDash \varphi(O_1,\dots,O_n,X)$. Since the extensions of $S_1',\dots,S_{2^n}'$ are all disjoint, to define $A$ it suffices to identify subsets of each $\semantics{S_k'}_{\mathcal{M}}$. As above, suppose $S'_k$ is equivalent to $S_i \vee S_j$, so that $S_i$ is equivalent to $S'_k \wedge X$ and $S_j$ is equivalent to $S_k' \wedge \neg X$. Then let $B_k$ be any subset of $\semantics{S_k'}_{\mathcal{M}}$ of size $b_i + a_{i,1}\mathsf{u}_1 + \dots + a_{i,m}\mathsf{u}_m$, such that the complement $\semantics{S'_k}_{\mathcal{M}} - B_k$ has size $b_j + a_{j,1}\mathsf{u}_1 + \dots + a_{j,m}\mathsf{u}_m$. This is always possible since $|S'_k|$ is simply the sum of these two numbers. Finally let $A = \bigcup_{k\leq 2^n} B_k$. 

Once again absorbing $X$ into the language and defining $\mathcal{M}'$ to be just like $\mathcal{M}$ but with $X^{\mathcal{M}'} = A$, the tuple $[|S_1|_{\mathcal{M}'},\dots,|S_{2^{n+1}}|_{\mathcal{M}'}]$ gives a solution to (\ref{semi-lin}) with the same choices $\mathsf{u}_1,\dots,\mathsf{u}_m$. Hence, $\mathcal{M}'\vDash \varphi(O_1,\dots,O_n,X)$, from which it easily follows that $\mathcal{M},s^X_A \vDash \varphi(O_1,\dots,O_n,X)$ and finally $\mathcal{M} \vDash \exists X. \varphi(O_1,\dots,O_n,X)$.
\end{proof}

\end{document}
